% \vspace{-.25cm}
\section{Related Work}
\label{sec:related}
% \vspace{-.25cm}

Recent advances in sequence modeling with deep networks have led to rapid improvement in the effectiveness of such models, from LSTMs and sequence-to-sequence models~\citep{hochreiter1997long, NIPS2014_a14ac55a} to Transformer architectures with self-attention \citep{vaswani2017attention}.
In light of this, it is tempting to consider how such sequence models can lead to improved performance in RL, which is also concerned with sequential processes~\citep{sutton1988learning}. Indeed, a number of prior works have studied applying sequence models of various types to represent components in standard RL algorithms, such as policies, value functions, and models~\citep{bakker2002reinforcement, Heess2015MemorybasedCW, chiappa2017recurrent,parisotto2020stabilizing,parisotto2021efficient,kumar2020adaptive}. While such works demonstrate the importance of such models for representing memory~\citep{oh2016control}, they still rely on standard RL algorithmic advances to improve performance. The goal in our work is different: we aim to replace as much of the RL pipeline as possible with sequence modeling, so as to produce a simpler method whose effectiveness is determined by the representational capacity of the sequence model rather than algorithmic sophistication.

Estimation of probability distributions and densities arises in many places in learning-based control. This is most obvious in model-based RL, where it is used to train predictive models that can then be used for planning or policy learning~\citep{sutton1990dyna,silver2008dyna2,fairbank2008reinforcement,deisenroth2011pilco,lampe2014modelnfq,heess2015svg,janner2020gamma,amos2020model}. However, it also figures heavily in offline RL, where it is used to estimate conditional distributions over actions that serve to constrain the learned policy to avoid out-of-distribution behavior that is not supported under the dataset~\citep{fujimoto2019off,kumar2019stabilizing,ghasemipour2020emaq}; imitation learning, where it is used to fit an expert's actions to obtain a policy~\citep{ross2010efficient,ross2011reduction}; and other areas such as hierarchical RL~\citep{peng2017deeploco,co2018self,jiang2019language}.
In our method, we train a single high-capacity sequence model to represent the joint distribution over sequences of states, actions, and rewards.
This serves as \emph{both} a predictive model \emph{and} a behavior policy (for imitation) or behavior constraint (for offline RL).

Our approach to RL is most closely related to prior model-based methods that plan with a learned model~\citep{chua2018pets,wang2019exploring}.
However, while these prior methods typically require additional machinery to work well, such as ensembles in the online setting~\citep{kurutach2018model,buckman2018sample,malik2019calibrated} or conservatism mechanisms in the offline setting~\citep{yu2020mopo,kidambi2020morel,argenson2020model}, our method does not require explicit handling of these components.
Modeling the states and actions jointly already provides a bias toward generating in-distribution actions, which avoids the need for explicit pessimism~\citep{fujimoto2019off,kumar2019stabilizing,ghasemipour2020emaq,nair2020accelerating,jin2020pessimism,yin2021near,dadashi2021offline}.
Our method also differs from most prior model-based algorithms in the dynamics model architecture used, with fully-connected networks parameterizing diagonal-covariance Gaussian distributions being a common choice \citep{chua2018pets}, though recent work has highlighted the effectiveness of autoregressive state prediction \citep{zhang2021autoregressive} like that used by the Trajectory Transformer.
In the context of recently proposed offline RL algorithms, our method can be interpreted as a combination of model-based RL and policy constraints~\citep{kumar2019stabilizing,wu2019behavior}, though our approach does not require introducing such constraints explicitly.
In the context of model-free RL, our method also resembles recently proposed work on goal relabeling \citep{andrychowicz2017hindsight,rauber2017hindsight,ghosh2019gcsl,paster2021planning} and reward conditioning \citep{schmidhuber2019reinforcement,srivastava2019training,kumar2019reward} to reinterpret all past experience as useful demonstrations with proper contextualization.

Concurrently with our work, \citet{chen2021decision} also proposed an RL approach centered around sequence prediction, focusing on reward conditioning as opposed to the beam-search-based planning used by the Trajectory Transformer.
Their work further supports the possibility that a high-capacity sequence model can be applied to reinforcement learning problems without the need for the components usually associated with RL algorithms.
